\section{Algebraic Integers}

\begin{definition}[Algebraic Integer]
    \label{def:algebraic-integer}%
    A complex number \(\alpha \in \CC\) is an \emph{algebraic integer} if it is the root of a monic polynomial in \(\ZZ\bqty{x}\).

    We define
    \[
        \ZZ\bqty{\alpha} = \bigcap_{\alpha \in R \subring \CC} R,
    \]
    or equivalently, for the \emph{evaluation homomorphism}
    \[
        \function{e_\alpha}{\ZZ\bqty{x}}{\CC}{g\pqty{x}}{g\pqty{\alpha},}
    \]
    we define
    \[
        \ZZ\bqty{\alpha} = \img e_{\alpha}.
    \]
\end{definition}

Note that \(\pi\) is not an algebraic integer, but nonetheless we may put \(\ZZ\bqty{\pi}\) into this definition. However, this process never `terminates', and this process is clearly very differently behaved for such transcendental numbers and algebraic integers.

\begin{proposition}[Ideal of Polynomial with an Algebraic Integer as Root is Principle]
    \label{prop:ideal-polynomial-algebraic-integer-principle}%
    Let \(\alpha\) be an algebraic integer. The ideal
    \[
        I = \ker e_{\alpha} \idealsub \ZZ\bqty{x}
    \]
    is \emph{principle}. It is equal to \(\pqty{f_\alpha}\) with \(f_\alpha\) irreducible and monic.
\end{proposition}

Note this is a non-trivial result! \(\ZZ\bqty{x}\) is not a PID.

\begin{proof}
    First, we note
    \[
        \ker e_\alpha = I
    \]
    is non-zero. Take \(f_\alpha \in I\) with non-trivial minimal degree. We can assume that \(f_\alpha\) is primitive, because \(\alpha\) satisfies a monic equation. We want to show \(I = \pqty{f_{\alpha}}\).

    Let \(h \in I\). In \(\QQ\bqty{x}\) we have division algorithm, so
    \[
        h = f_{\alpha} \cdot q + r
    \]
    with \(\deg\pqty{r} < \deg\pqty{f_{\alpha}}\). We clear the denominators, such that
    \[
        ah = af_{\alpha} \cdot q + ar
    \]
    for \(a \in \ZZ\). Putting in \(\alpha\), we have \(\pqty{ah} \pqty{\alpha} = 0\), and \(f_{\alpha} \pqty{\alpha} = 0\), so \(\pqty{ar} \pqty{\alpha} = 0\). By the minimality of the degree of \(f_\alpha\), we get \(\pqty{ar} = 0\) is the zero polynomial, so \(r = 0\).

    Therefore, \(ah = af_{\alpha} \cdot q\). But recall \(f_{\alpha}\) is primitive. Taking contents, we have
    \[
        \content\pqty{ah} = \content\pqty{aq}.
    \]
    
    We see that \(h \in I \idealsub \ZZ\bqty{x}\), and \(a \divides \content\pqty{ah}\), so \(a \divides \content\pqty{aq}\). Therefore,
    \[
        aq = a \bar{q}
    \]
    so \(\bar{q} \in \ZZ\bqty{x}\). Cancelling \(a\)s gives \(q = \bar{q}\), so \(f_\alpha\) divides \(h\) as claimed, i.e.,
    \[
        I = \pqty{f_{\alpha}}.
    \]

    It remains to show that \(f_{\alpha}\) is irreducible, note that
    \[
        \flatfrac{\ZZ\bqty{x}}{\ker\pqty{e_\alpha}} \isomorphic \flatfrac{\ZZ\bqty{x}}{\pqty{f_{\alpha}}} \subring \CC
    \]
    so \(\flatfrac{\ZZ\bqty{x}}{\pqty{f_{\alpha}}}\) is an integral domain. It follows that \(f_\alpha\) is prime, hence irreducible.
\end{proof}

\begin{example}[Algebraic Integers]
    \label{ex:algebraic-int}%
    Note that \(\sqrt{2}\), \(i\), \(\sqrt{-3}\) are algebraic integers. Furthermore, 
    \[
        \frac{1 + \sqrt{-3}}{2}
    \]
    is an algebraic integer, since it is a root of unity.
\end{example}

\begin{example}[Non-Algebraic Integers]
    \label{ex:non-algebraic-int}%
    \(\pi\), \(e\) are examples of non-algebraic integers. In fact, algebraic integers form a ring, but this is much harder to prove.
\end{example}

\begin{proposition}[Rational Algebraic Integers are Integers]
    \label{prop:rational-algebraic-integer}%
    Let \(\alpha \in \CC\) be an algebraic integer. If \(\alpha \in \QQ\), then \(\alpha \in \ZZ\).
\end{proposition}

\begin{proof}
    Let \(f_\alpha\) be the minimal polynomial generating \(\ker e_\alpha\). This is irreducible over \(\ZZ\bqty{x}\), so also \(\QQ\bqty{x}\) by \autoref{thm:gauss-lemma} (Gauss's Lemma). But \(\pqty{x - \alpha}\) divides \(f_\alpha\) in \(\QQ\bqty{x}\), so \(f_\alpha = x - \alpha\).
\end{proof}